% ============== Section 5: Reconfigurable Intelligent Surfaces (S30-S35) ==
\section{Reconfigurable Intelligent Surfaces}

% --- S30 -----------------------------------------------------------------
\begin{frame}{RIS: programming the channel itself}
  \centering
  \resizebox{0.8\textwidth}{!}{% RIS-aided link geometry (slides 30/31) -- native TikZ
\begin{tikzpicture}
  % BS
  \node[blkb, minimum width=13mm] (bs) at (0,0) {BS};
  \draw[thick] (bs.north) -- ++(0,3mm) -- ++(-2mm,3mm) ++(2mm,-3mm) -- ++(2mm,3mm);
  % RIS panel
  \begin{scope}[shift={(5.2,2.3)}, rotate=-14]
    \draw[thick, fill=accent!12] (-1.15,-0.62) rectangle (1.15,0.62);
    \foreach \x in {-0.92,-0.46,0,0.46,0.92}
      \foreach \y in {-0.4,0,0.4}
        \draw[fill=accent!45, draw=accent] (\x-0.14,\y-0.11) rectangle (\x+0.14,\y+0.11);
  \end{scope}
  \node[lab] at (5.2,3.45) {RIS: $N$ tunable elements,
    $\Theta=\mathrm{diag}(\beta_n e^{j\theta_n})$};
  % UE
  \node[blk, minimum width=11mm] (ue) at (9.6,-0.2) {UE};
  % obstacle
  \draw[fill=gray!35, draw=gray] (4.35,-0.85) rectangle (5.35,0.75);
  \node[lab, gray] at (4.85,-1.12) {blockage};
  % links
  \draw[arr, cblue, very thick] (0.8,0.55) -- node[lab, above, sloped]{$G$ (BS$\to$RIS)} (4.25,2.15);
  \draw[arr, cblue, very thick] (6.2,2.1) -- node[lab, above, sloped]{$h_r$ (RIS$\to$UE)} (9.25,0.35);
  \draw[arr, gray, dashed, thick] (0.9,0.0) -- node[lab, below]{$h_d$ (weak / blocked)} (9.0,-0.2);
  % equation
  \node[lab, align=center] at (4.9,-1.9)
    {$y=\left(\bm h_r^{H}\,\Theta\, \bm G + h_d^{H}\right)x+n
     \qquad\text{path loss}\ \propto (d_1 d_2)^{-\alpha}\ \text{(dyadic)}$};
\end{tikzpicture}
}
  \note{[S30, 1 min | clock 0:40, after Demo 2] Paradigm framing: for a
  century the channel was fate --- estimate it, adapt to it. A RIS makes part
  of the environment a design variable: hundreds of passive elements, each
  adding a controllable phase to the reflection, steered by a low-rate control
  link. No PAs, no ADCs on the panel.}
\end{frame}

% --- S31 -----------------------------------------------------------------
\begin{frame}{The cascaded channel model}
  \begin{equation*}
    y=\left(\bm h_r^{H}\,\Theta\,\bm G+h_d^{H}\right)x+n,
    \qquad
    \Theta=\mathrm{diag}\!\left(\beta_1 e^{j\theta_1},\dots,
                                \beta_N e^{j\theta_N}\right)
  \end{equation*}
  \begin{itemize}\itemsep2pt
    \item Product path loss $\propto (d_1 d_2)^{-\alpha}$: the
          \alert{fundamental tax of passivity} --- both hops paid in full
    \item Estimation burden: cascaded coefficients scale as $N\!\times\!M$;
          structure (sparsity, common subspaces) to the rescue
    \item Large $N$ at short range $\Rightarrow$ the RIS itself operates in the
          near field
  \end{itemize}
  \note{[S31, 1 min | clock 0:41] The equation is deceptively simple; the two
  hard truths live under it: dyadic path loss (why RIS placement matters
  enormously) and the estimation burden (you cannot observe the RIS channel
  directly --- the panel has no receivers). Both are active research veins.}
\end{frame}

% --- S32 -----------------------------------------------------------------
\begin{frame}{Phase optimization and the $N^2$ law}
  \begin{equation*}
    \theta_n^{\star}=\arg(h_d)-\arg\!\left([\bm h_r]_n^{*}[\bm G]_n\right)
    \;\;\Rightarrow\;\;
    P_{\mathrm{rx}}\propto\Big(\textstyle\sum_n |[\bm h_r]_n||[\bm G]_n|\Big)^{2}
    \sim N^{2}
  \end{equation*}
  \begin{columns}[T]
    \column{0.5\textwidth}
    \begin{itemize}\itemsep1pt
      \item Co-phasing: amplitude $\propto N$ $\Rightarrow$ power
            $\propto N^2$ ($+6$\,dB per doubling)
      \item Random phases: 2D random walk $\Rightarrow$ power $\propto N$
      \item Multi-user / MIMO: non-convex $\Rightarrow$ SDR, manifold,
            alternating optimization
    \end{itemize}
    \column{0.48\textwidth}
    \centering
    \includegraphics[width=\textwidth]{fig_demo3_ris_scaling.png}
  \end{columns}
  \note{[S32, 1.5 min | clock 0:42] Derive the closed form in one breath ---
  single-user SNR max separates per element; every element should cancel its
  cascade phase. The plot is our simulation with the analytic overlays: slopes
  2.0 and 1.0 on log-log. But remember the $N^2$ rides on top of the dyadic
  loss --- this is Demo 3's punchline in three minutes.}
\end{frame}

% --- S33 -----------------------------------------------------------------
\begin{frame}{Active RIS and other variants}
  \begin{equation*}
    y=\left(\bm h_r^{H}\Theta\,\bm G\right)x
      +\underbrace{\bm h_r^{H}\Theta\,\bm n_{\mathrm{RIS}}}_{\text{amplified noise}}
      +n
  \end{equation*}
  \begin{itemize}\itemsep2pt
    \item Passive gain often eaten by dyadic loss $\Rightarrow$ \alert{active
          RIS}: reflection with amplification --- and an amplified-noise tax
    \item Variants: STAR-RIS (transmit \emph{and} reflect), beyond-diagonal
          RIS, holographic surfaces
    \item The uncomfortable benchmark: does an active RIS beat a plain relay /
          small cell at the same cost and power? Must be argued per scenario
  \end{itemize}
  \note{[S33, 1 min | clock 0:43:30] Keep the field honest: the moment you add
  amplifiers you must compare against the century-old alternative --- a relay.
  The interesting regime for passive RIS: cheap, dense, conformal surfaces
  where relays are impractical (facades, indoor panels). STAR-RIS one-liner:
  serve both sides of the surface.}
\end{frame}

% --- S34 -----------------------------------------------------------------
\begin{frame}{Deployment: where RIS actually helps}
  \begin{itemize}\itemsep2pt
    \item Sweet spots: mmWave/sub-THz \alert{blockage bypass}, indoor dead
          zones, localized capacity boosts
    \item Placement rule of thumb: near TX or near RX --- mid-path maximizes the
          dyadic loss penalty
    \item Overheads that gate practicality: channel estimation and the control
          link, both scaling with $N$
    \item Field trials (metro, industrial) encouraging; standardization
          pre-normative (ETSI ISG RIS; 3GPP interest)
  \end{itemize}
  \note{[S34, 1 min | clock 0:44:30] Deployment intuition via the dyadic law:
  $(d_1 d_2)$ is minimized at the ends, maximized in the middle --- put the
  mirror next to the lamp or next to the book, not halfway. The overhead
  bullet is the researcher's opening: estimation with almost no observations
  is a beautiful problem.}
\end{frame}

% --- S35 -----------------------------------------------------------------
\begin{frame}{RIS: open problems}
  \begin{itemize}\itemsep2pt
    \item CSI acquisition with no RF chains on the panel
    \item Wideband behavior: one phase per element, many subcarriers
    \item EM-consistent models: mutual coupling, non-ideal reflection
    \item RIS-aided sensing and localization: surfaces as programmable
          landmarks
  \end{itemize}
  \demoslot{3 (2.5 min)}{\textbf{The $N^2$ law, live}: SNR gain from 16 to
    1024 elements, random vs optimized phases.
    \hfill\src{matlab/demo3\_ris\_gain}}
  \note{[S35 + DEMO 3, 0.5+2.5 min | clock 0:45:30] Rapid-fire open list, then
  MATLAB. Demo 3 is the shortest: the two scaling lines draw themselves; ask
  the audience to call the slope difference before the fit prints (2 vs 1).
  If MATLAB fails: appendix A3. Demo done by ~0:48:30.}
\end{frame}
